\begin{frame}{Änderungsanträge (LIBE)}
\tiny
\vspace{-0.45cm}
\renewcommand{\arraystretch}{0.9}
\setlength{\tabcolsep}{2pt}
\setlength{\abovetopsep}{0pt}
\setlength{\belowrulesep}{0pt}
\begin{tabularx}{\textwidth}{@{}>{\centering\arraybackslash}m{1.55cm} >{\RaggedRight\arraybackslash}m{\dimexpr\linewidth-1.55cm-2\tabcolsep\relax}@{}}
\toprule
\multicolumn{1}{m{1.55cm}}{\centering\textbf{Fraktion}} & \multicolumn{1}{m{\dimexpr\linewidth-1.55cm-2\tabcolsep\relax}}{\centering\textbf{Änderungsantrag}} \\
\midrule
\includegraphics[width=2cm,height=1cm,keepaspectratio]{Bilder/EVP.png} & 11. ... regt an, Regelungen für eine verpflichtende, einjährige Wehrpflicht bei den Streitkräften der Europäischen Union \textcolor{red}{\sout{im ersten Jahr}}\textcolor{green}{\uline{nach Vollendung}} der \textcolor{red}{\sout{Volljährigkeit}}\textcolor{green}{\uline{aktüllen Ausbildung und Vollendung des 18. Lebensjahres}} für alle Unionsbürgerinnen und -bürger zu schaffen. Somit wird die Wehrfähigkeit der \\
\includegraphics[width=2cm,height=1cm,keepaspectratio]{Bilder/SD.png} & 11. ... regt an, Regelungen für eine \textcolor{red}{\sout{verpflichtende, einjährige}}\textcolor{green}{\uline{freiwillige, mindestens sechsmonatige}} Wehrpflicht bei den Streitkräften der Europäischen Union \textcolor{red}{\sout{im ersten Jahr}}\textcolor{green}{\uline{nach der Erreichung}} der Volljährigkeit \textcolor{green}{\uline{bis zum 40. Lebensjahr }}für alle Unionsbürgerinnen und -bürger zu schaffen. Somit wird die Wehrfähigkeit der \\
\includegraphics[width=2cm,height=1cm,keepaspectratio]{Bilder/RE.png} & Union sichergestellt. Das Recht, den Dienst an der Waffe zu verweigern, ist dabei jedoch unbedingt anzürkennen.\textcolor{green}{\uline{ 12. Um die persönliche Freiheit der Bürgerinnen und Bürger zu gewährleisten und Ausgaben möglichst effizient zu gestalten, sind wir gegen eine Wehrpflicht. Wir wollen den Bürgerinnen und Bürgern die Entscheidung lassen, freiwillig Wehrdienst zu leisten oder sich in einem sozialen Dienst zu engagieren.}} \\
\includegraphics[width=2cm,height=1cm,keepaspectratio]{Bilder/PfE.png} & \textcolor{red}{\sout{11. ... regt an, Regelungen für eine verpflichtende, einjährige Wehrpflicht bei den Streitkräften der Europäischen Union im ersten Jahr der Volljährigkeit für alle Unionsbürgerinnen und -bürger zu schaffen.}}\textcolor{green}{\uline{11. ... regt an, ein Konzept zu entwickeln für eine einjährige Wehrpflicht bei den nationalen Streitkräften d im ersten Jahr der Volljährigkeit, das von den Mitgliedsstaaten übernommen und angepasst werden kann.}} Somit wird die Wehrfähigkeit der Union sichergestellt. Das Recht, den Dienst an der Waffe zu verweigern, ist \\
\includegraphics[width=2cm,height=1cm,keepaspectratio]{Bilder/Left.png} & 11. ... \textcolor{red}{\sout{regt an, Regelungen für eine verpflichtende, einjährige Wehrpflicht bei den Streitkräften der Europäischen Union im ersten Jahr der Volljährigkeit für alle Unionsbürgerinnen und -bürger zu schaffen. Somit wird die Wehrfähigkeit der Union sichergestellt. Das Recht, den Dienst an der Waffe zu verweigern, ist dabei jedoch unbedingt anzürkennen.}}\textcolor{green}{\uline{regt an, Regelungen für einen optionalen und freiwilligen Wehrdienst bei den Streitkräften der Europäischen Union ab der Vollendung des 18. Lebensjahrs für alle Unionsbürgerinnen und -bürger zu schaffen. Somit wird die Wehrfähigkeit der Union sichergestellt und gleichzeitig die Optionsfreiheit des Einzelnen bewahrt.}} \\
\bottomrule
\end{tabularx}
\end{frame}

